% !PRINTED 351
\weberpage{351}

To obtain the imaginary roots, one replaces $\mathrm{tg}\,\varphi$ by $\varrho\,\mathrm{tg}\,\varphi$, where $\varrho$ denotes an imaginary third root of unity.

b) In the second case we determine the angle $\vartheta$, likewise in the first quadrant, from

\weq{7}{\frac{g}{2} = \sqrt{\frac{e^3}{27}}\,\frac{1}{\sin\vartheta}}

\[
\sqrt{R} = \sqrt{\frac{e^3}{27}}\,\mathrm{cotg}\,\vartheta
\]

\[
x = \sqrt{\frac{e}{3}}\left(\sqrt[3]{\mathrm{cotg}\,\frac{\vartheta}{2}} + \sqrt[3]{\mathrm{tg}\,\frac{\vartheta}{2}}\right),
\]

\weq{8}{\mathrm{tg}\,\frac{\vartheta}{2} = \mathrm{tg}\,\varphi^3,}

\weq{9}{x = 2\sqrt{\frac{e}{3}}\,\frac{1}{\sin 2\varphi}.}

c) In the last case, finally, where three real roots are present, we set

\weq{10}{\frac{g}{2} = \sqrt{\frac{e^3}{27}}\cos\vartheta}

\[
\sqrt{R} = i\sqrt{\frac{e^3}{27}}\sin\vartheta
\]

\[
x = \sqrt{\frac{e}{3}}\left(\sqrt[3]{\cos\vartheta + i\sin\vartheta} + \sqrt[3]{\cos\vartheta - i\sin\vartheta}\right),
\]

or, by the theorem of \gesperrt{Moivre}:

\weq{11}{x = 2\sqrt{\frac{e}{3}}\cos\frac{1}{3}\vartheta.}

If one takes $\vartheta$ again in the first quadrant, one obtains for the two other, likewise real, roots

\[
-2\sqrt{\frac{e}{3}}\cos\frac{\pi + \vartheta}{3}, \qquad -2\sqrt{\frac{e}{3}}\cos\frac{\pi - \vartheta}{3}.
\]

All these formulae are arranged for logarithmic computation.
